% Einheit 2 von 4 – Scheitelpunktform (bank/quadratische-funktionen/e2.jsonl, zusammenbau v0.1)
%% TODO Zweigzeile Teil 1: was hier gelernt wird (Satz in Schülersprache aus dem Einheitstitel) – Einheit 2
\einheitenkopf[e2]{Einheit 2 von 4 · Scheitelpunktform}
\zweigzeile{ab Kl. 9, je nach Buch bis Kl. 10 · P10 oft}
\setcounter{aufgabe}{12}

%% TODO Titel: Ich-kann-Satz fehlt in der Bank (Platzhalter: Kettenname) – Nr. 13
%% TODO Anweisung: ein Satz über der ersten Teilaufgabe fehlt in der Bank – Nr. 13
\begin{aufgabe}{Scheitelpunktform}
\begin{teile}
\teil $f(x) = (x - 4)^2 + 2$ – kreise die Zahl in der Klammer ein. Bei welchem $x$ liegt der Scheitel? x = \leerfeld
\end{teile}
%% TODO Darstellung für schwach fehlt in der Bank (quadratische-funktionen-e2-k1-s1-v1; Abschnitt „Für schwache Schüler“)
\swz{$f(x) = (x - 2)^2 + 5$ – Scheitel?}{}
%% TODO Darstellung für schwach fehlt in der Bank (quadratische-funktionen-e2-k1-s1-v2; Abschnitt „Für schwache Schüler“)
\swz{$f(x) = (x - 4)^2 + 1$ – Scheitel?}{}
%% TODO Darstellung für schwach fehlt in der Bank (quadratische-funktionen-e2-k1-s1-v3; Abschnitt „Für schwache Schüler“)
\swz{$f(x) = (x - 5)^2 + 3$ – Scheitel?}{}
%% TODO Darstellung für schwach fehlt in der Bank (quadratische-funktionen-e2-k1-s1-v4; Abschnitt „Für schwache Schüler“)
\swz{$f(x) = (x - 1)^2 + 6$ – Scheitel?}{}
%% TODO Darstellung für schwach fehlt in der Bank (quadratische-funktionen-e2-k1-s1-v5; Abschnitt „Für schwache Schüler“)
\swz{$f(x) = (x - 6)^2 + 2$ – Scheitel?}{}
\swfrage{Was bleibt gleich, was ändert sich?}
%% TODO Darstellung für schwach fehlt in der Bank (quadratische-funktionen-e2-k1-s2-v1; Abschnitt „Für schwache Schüler“)
\swz{$f(x) = (x + 4)^2 + 2$ – Scheitel?}{}
\swz{Lies den Scheitel der Parabel ab.}{\begin{ksys}[xmin=-5,xmax=1,ymin=-1,ymax=8,ablesen]
\parabel{1}{-2}{3}{f}
\end{ksys}}
\swa{$f(x) = (x - 2)^2 + 1$ – skizziere die Parabel.}{\begin{ksys}[xmin=-2,xmax=6,ymin=-1,ymax=10]
\end{ksys}}
%% TODO Darstellung für schwach fehlt in der Bank (quadratische-funktionen-e2-k1-s5-v1; Abschnitt „Für schwache Schüler“)
\swz{Nach oben geöffnete Normalparabel mit Scheitel $S(5|2)$ – Gleichung?}{}
\swz{Gleichung der Parabel?}{\begin{ksys}[xmin=-1,xmax=5,ymin=-4,ymax=2,ablesen]
\parabel{1}{2}{-3}{f}
\end{ksys}}
%% TODO Darstellung für schwach fehlt in der Bank (quadratische-funktionen-e2-k1-s7-v1; Abschnitt „Für schwache Schüler“)
\swz{$f(x) = (x - 1)^2 + 2$ wird um $4$ nach unten verschoben – Gleichung der neuen Parabel $g$?}{}
%% TODO Darstellung für schwach fehlt in der Bank (quadratische-funktionen-e2-k1-s8-v1; Abschnitt „Für schwache Schüler“)
\swz{$f(x) = (x - 5)^2 + 1$ wird an der $x$-Achse gespiegelt – Gleichung des Spiegelbilds $g$?}{}
%% TODO Darstellung für schwach fehlt in der Bank (quadratische-funktionen-e2-k1-s9-v1; Abschnitt „Für schwache Schüler“)
\swz{$f(x) = (x - 4)^2 + 2$ wird an der $y$-Achse gespiegelt – Gleichung des Spiegelbilds $g$?}{}
%% TODO Darstellung für schwach fehlt in der Bank (quadratische-funktionen-e2-k1-s10-v1; Abschnitt „Für schwache Schüler“)
\swz{Gib eine nach oben geöffnete Normalparabel an, deren Scheitel auf der $x$-Achse rechts vom Ursprung liegt.}{}
%% TODO Darstellung für schwach fehlt in der Bank (quadratische-funktionen-e2-k1-s11-v1; Abschnitt „Für schwache Schüler“)
\swz[3]{$f(x) = (x - 1)^2 + 2$ und $g(x) = (x - 1)^2 - 4$ – wie liegen die Parabeln zueinander? Haben sie gemeinsame Punkte? Begründe ohne Rechnung.}{}
%% TODO Darstellung für schwach fehlt in der Bank (quadratische-funktionen-e2-k1-s12-v1; Abschnitt „Für schwache Schüler“)
\swz{$f(x) = 2(x - 4)^2 + 1$ – Scheitel, Öffnung, schmaler oder breiter als die Normalparabel? Punkt eine Einheit rechts vom Scheitel?}{}
\swz[3]{$f(x) = (x - 4)^2 - 5$ wird erst um $5$ nach oben verschoben und dann an der $x$-Achse gespiegelt. Gleichung der neuen Parabel $g$? \hfill (P10 2017 OS)}{}
\swa{Die Parabel $f(x) = (x + 4)^2 - 2$ wird an der $y$-Achse gespiegelt. Skizziere das Spiegelbild $g$ und gib seine Gleichung an. \hfill (P10 2020 OS) g(x) = \leerfeld}{\begin{ksys}[xmin=-7,xmax=7,ymin=-3,ymax=7]
\parabel{1}{-4}{-2}{f}
\end{ksys}}
\swa{Skizziere eine nach unten geöffnete Normalparabel $f$ mit Scheitel auf der $y$-Achse, die keine Nullstelle hat. Gib ihre Gleichung und die Gleichung ihres Spiegelbilds $g$ an der $x$-Achse an. \hfill (P10 2018 OS) f(x) = \leerfeld , g(x) = \leerfeld}{\begin{ksys}[xmin=-4,xmax=4,ymin=-10,ymax=10,ystep=2]
\end{ksys}}
\end{aufgabe}

%% TODO Titel: Ich-kann-Satz fehlt in der Bank (Platzhalter: Kettenname) – Nr. 14
%% TODO Anweisung: ein Satz über der ersten Teilaufgabe fehlt in der Bank – Nr. 14
\begin{aufgabe}{Scheitelpunktform – Fehler finden, Begründen, Darstellungswechsel, Anwendung}
\swz[3]{Ben liest bei $f(x) = (x + 6)^2 - 2$ den Scheitel $S(6|{-2})$ ab. Finde den Fehler und korrigiere.}{}
\swfrage{Warum hat $f(x) = (x - 4)^2 + 1$ bei $x = 4$ ihren kleinsten Wert?}
\begin{teile}
\teil Wahr oder falsch? \\ Die Gleichung ist $f(x) = (x + 1)^2 - 4$. \kreuz{wahr} \kreuz{falsch} \\ Der Scheitel ist $S(1|{-4})$. \kreuz{wahr} \kreuz{falsch} \\ Die Parabel ist nach unten geöffnet. \kreuz{wahr} \kreuz{falsch}
\end{teile}
\swz[3]{Beim Kugelstoßen gilt für die Flugbahn $h(x) = -0{,}05(x - 6)^2 + 3{,}8$ ($x$: Entfernung in m, $h$: Höhe in m). Wie hoch fliegt die Kugel höchstens, und in welcher Entfernung?}{\begin{ksys}[xmin=0,xmax=16,ymin=0,ymax=5,xstep=2,xlabel=$x$ in m,ylabel=$h$ in m]
\funktionab{-0.05*(\x-6)^2+3.8}{h}{0}{14.7}
\end{ksys}}
\end{aufgabe}

\uebersichtskasten{Scheitelpunktform: p(x) = (x − d)² + e hat den Scheitel S(d | e) – in der Klammer steht d mit umgedrehtem Vorzeichen. \\ (x − 3)² + 1 → S(3 | 1) \quad (x + 3)² − 1 → S(−3 | −1) \quad −(x − 3)² + 1 → S(3 | 1), nach unten geöffnet \\ Zeichnen: Scheitel setzen, dann wie bei der Normalparabel weiter: 1 nach rechts 1 nach oben, 2 nach rechts 4 nach oben (Schablone). \\ Verschieben ändert nur d und e. Spiegeln an der x-Achse: Minus vor die Klammer und e umgedreht (−(x − 3)² − 1). Spiegeln an der y-Achse: d umgedreht ((x + 3)² + 1). \\ Nullstellen ohne Rechnung, nach oben geöffnet: Scheitel unter der x-Achse → zwei, auf der x-Achse → eine, darüber → keine.}
